\documentclass[11pt,a4paper]{article}
\usepackage[utf8]{inputenc}
\usepackage[spanish]{babel}
\usepackage{listings}
\usepackage{xcolor}
\usepackage{hyperref}
\usepackage{geometry}
\geometry{margin=1in}

% Configuración de colores para código fuente
\definecolor{codegreen}{rgb}{0,0.6,0}
\definecolor{codegray}{rgb}{0.5,0.5,0.5}
\definecolor{codepurple}{rgb}{0.58,0,0.82}
\definecolor{backcolour}{rgb}{0.08,0.11,0.17}
\definecolor{codefg}{rgb}{0.9,0.9,0.9}

\lstdefinestyle{mystyle}{
    backgroundcolor=\color{backcolour},   
    commentstyle=\color{codegreen},
    keywordstyle=\color{orange},
    numberstyle=\tiny\color{codegray},
    stringstyle=\color{codepurple},
    basicstyle=\ttfamily\footnotesize\color{codefg},
    breakatwhitespace=false,         
    breaklines=true,                 
    captionpos=b,                    
    keepspaces=true,                 
    numbers=left,                    
    numbersep=5pt,                  
    showspaces=false,                
    showstringspaces=false,
    showtabs=false,                  
    tabsize=2
}

\lstset{style=mystyle}

\title{Manual Técnico: Sistema de Administración, Seguridad y Auditoría}
\author{AEGIS Wiki - Elden Ring}
\date{\today}

\begin{document}

\maketitle

\section{Introducción}
Este manual detalla la arquitectura, el diseño de la base de datos, los endpoints de la API, las políticas de seguridad y el sistema de auditoría implementados para cumplir con los requerimientos técnicos de las \textbf{Prácticas 11 y 12}.

\section{Arquitectura General del Sistema}
El sistema sigue una arquitectura desacoplada cliente-servidor (SPA + API REST):

\begin{itemize}
    \item \textbf{Frontend}: Angular 21 (SPA) con componentes standalone y signals para reactividad. Utiliza PrimeNG como biblioteca de componentes de interfaz.
    \item \textbf{Backend}: Node.js con Express, arquitectura REST y pool de conexiones PostgreSQL.
    \item \textbf{Base de Datos}: PostgreSQL 17 hospedado en la nube mediante \textbf{Neon}, con cifrado SSL obligatorio y conexión pooled.
    \item \textbf{Servidor Backend}: Desplegado en \textbf{Render} dentro de un contenedor Docker automatizado.
    \item \textbf{Servidor Frontend}: Desplegado en \textbf{Vercel} con reescritura de API dinámica (\texttt{vercel.json}).
\end{itemize}

\section{Estructura de la Base de Datos}
Las tablas involucradas en el control de acceso, seguridad y auditoría en Neon PostgreSQL son:

\subsection{Tabla: \texttt{usuarios}}
Almacena las cuentas de usuario y los estados de bloqueo/seguridad:
\begin{lstlisting}[language=SQL, caption=Definición física de la tabla de usuarios]
CREATE TABLE IF NOT EXISTS usuarios (
  id SERIAL PRIMARY KEY,
  nombre VARCHAR(80) NOT NULL,
  email VARCHAR(120) UNIQUE NOT NULL,
  role VARCHAR(16) NOT NULL DEFAULT 'user' CHECK (role IN ('admin', 'user')),
  password_hash TEXT NOT NULL,
  activo BOOLEAN NOT NULL DEFAULT TRUE,                
  intentos_fallidos INT NOT NULL DEFAULT 0,            
  bloqueado_hasta TIMESTAMP DEFAULT NULL,              
  permisos TEXT[] NOT NULL DEFAULT '{}',               
  created_at TIMESTAMP NOT NULL DEFAULT NOW()
);
\end{lstlisting}

\subsection{Tabla: \texttt{bitacora} (Auditoría)}
Registra las acciones clave de forma inmutable:
\begin{lstlisting}[language=SQL, caption=Definición física de la tabla de bitácora]
CREATE TABLE IF NOT EXISTS bitacora (
  id SERIAL PRIMARY KEY,
  usuario VARCHAR(120) NOT NULL,                       
  fecha DATE NOT NULL DEFAULT CURRENT_DATE,            
  hora TIME NOT NULL DEFAULT CURRENT_TIME,            
  ip_address VARCHAR(45) NOT NULL,                     
  accion VARCHAR(255) NOT NULL                         
);
\end{lstlisting}

\section{Políticas y Controles de Seguridad Implementados}
El sistema cuenta con los siguientes controles de seguridad activa:

\begin{enumerate}
    \item \textbf{Cifrado Hash de Contraseñas}: Las claves se encriptan de manera unidireccional en el backend usando la biblioteca \texttt{bcrypt} con un costo de cómputo de 10 salt rounds (\texttt{bcrypt.hash(password, 10)}).
    \item \textbf{Políticas de Contraseña Segura}: Las contraseñas se evalúan para verificar que cumplan con una longitud mínima de 8 caracteres, al menos una letra mayúscula, una minúscula, un número y un carácter especial.
    \item \textbf{Bloqueo de Cuenta por Fuerza Bruta (Lockout)}: Si un usuario acumula 3 intentos fallidos de inicio de sesión de forma consecutiva, la columna \texttt{bloqueado\_hasta} se establece a \texttt{NOW() + 15 minutos}. El servidor deniega todas las peticiones posteriores hasta que se cumpla el plazo.
    \item \textbf{Baja Lógica (Soft Delete)}: Las cuentas eliminadas por el administrador no se borran físicamente mediante \texttt{DELETE} (lo cual provocaría huérfanos en auditoría). En su lugar, se actualiza la columna \texttt{activo} a \texttt{false} (\texttt{UPDATE usuarios SET activo = false}).
    \item \textbf{Expiración de Sesión Local}: El frontend almacena la sesión localmente y la hace expirar de forma transparente tras 30 minutos de inactividad, forzando un cierre de sesión.
\end{enumerate}

\section{Endpoints de la API (REST)}

\subsection{Endpoints de Autenticación y Perfil}
\begin{itemize}
    \item \texttt{POST /auth/register}: Registro de usuario aplicando políticas de contraseña.
    \item \texttt{POST /auth/login}: Validación de credenciales, control de bloqueo e inicio de sesión.
    \item \texttt{POST /auth/logout}: Registra la salida en la bitácora de auditoría.
    \item \texttt{PUT /auth/profile}: Edición del nombre del usuario de la sesión actual (auditable).
    \item \texttt{PUT /auth/password}: Cambio de contraseña propia (auditable).
\end{itemize}

\subsection{Endpoints de Administración (\texttt{x-user-role: admin} requerido)}
\begin{itemize}
    \item \texttt{GET /admin/usuarios}: Listar cuentas y permisos de todos los usuarios.
    \item \texttt{POST /admin/usuarios}: Crear un usuario asignando roles y permisos.
    \item \texttt{PUT /admin/usuarios/:id}: Actualizar datos, roles, permisos y estado activo de un usuario.
    \item \texttt{PUT /admin/usuarios/:id/password}: Reestablecer la clave de un usuario desde la administración.
    \item \texttt{DELETE /admin/usuarios/:id}: Desactivación lógica (baja) del usuario.
    \item \texttt{GET /admin/bitacora}: Obtener los registros de la bitácora de auditoría.
\end{itemize}

\end{document}
